% !TeX encoding = UTF-8
% !TeX spellcheck = sk_SK
% !TeX root=tukedip.tex
\setcounter{page}{1}
\setcounter{equation}{0}
\setcounter{figure}{0}
\setcounter{table}{0}

\section*{Úvod}
\addcontentsline{toc}{section}{\numberline{}Úvod}
Kľúčovú úlohu zohráva téma „Multimodálna detekcia príspevkov sociálnych sietí s obsahom vygenerovaných GenAI“, konkrétne rozpoznávanie umelej inteligencie na sociálnych sieťach, ako sú Instagram, Facebook a ďalšie. V súčasnosti je rozlišovanie medzi skutočným príspevkom a príspevkom vytvoreným umelou inteligenciou veľmi zložitou otázkou, na ktorú človek nedokáže na prvý pohľad okamžite odpovedať. Patrí sem aj určenie, ktorý príspevok sa práve nachádza pred ním, najmä preto, že nie všetci autori zanechajú poznámku, ktorá naznačuje, že daný príspevok bol vytvorený pomocou umelej inteligencie. 

Moja bakalárska práca obsahuje model, ktorý pomôže používateľovi identifikovať, ktorý príspevok sa nachádza pred ním, ako aj užívateľsky prívetivé rozšírenie, ktoré môže pomocou môjho modelu skontrolovať ktokoľvek bez znalostí programovania. 

Taktiež som vykonal výskum, v ktorom som testoval aspekty, ako je vplyv textu vytvoreného umelou inteligenciou na odpoveď modelu, ako aj obrázka vytvoreného umelou inteligenciou na odpoveď modelu, aby som zistil, ktorá časť má najväčší vplyv na odpoveď modelu.